\headline{{\textbf{\Large\sffamily Seasonal Care Tips:}}\\
          {\textbf{\large\sffamily Mid--February to Mid--March}}}
\phantomsection
\label{2026-02-care}
\vfill

\noindent
As the days perceptibly lengthen and the sun begins to hold a marginal degree of warmth, our collections stand on the precipice of the growing season.
The deep dormancy of January has shifted into a restless transitional phase.
While the thermometer may still hover around a modest 7°C to 10°C during the day, the biological clock of our trees is accelerating.
For the bonsai enthusiast, this is arguably the busiest month of the calendar.
The focus shifts abruptly from passive protection to active intervention, as we manage the narrow window between the first swelling of buds and the ``point of no return'' that is leaf\--burst.

\medskip
\topic{The Stirring of the Sap}
\noindent
Observe your trees daily; the transition can be remarkably swift.
\textbf{Larch (Larix)}, \textbf{Hawthorn}, and \textbf{Acer palmatum} are often the first to show movement, with buds changing from dry scales to vibrant, plump ``candles.''
Once the buds begin to swell, the tree is no longer fully dormant.
This means its frost hardiness is significantly reduced.
While a tree could handle −5°C in December, the same temperature now could kill emerging tissue or cause ``die-back'' on refined branches.
Be prepared to move early-starters back under glass or into a sheltered porch if a late-winter cold snap is forecast.

\medskip
\topic{The Repotting Marathon}
\noindent
This is the primary task for the coming four weeks.
The goal is to repot just as the buds swell but before the leaves unfurl.
Prioritize \textbf{deciduous species} first, followed by \textbf{conifers} and finally \textbf{evergreens}.
When root pruning, be clinical: remove thick, downward-growing ``anchor'' roots to encourage a flat, fibrous \textbf{nebari}.
Ensure your \textbf{Akadama, Pumice, and Kiryu} mixes are bone\--dry before use; using wet substrate during repotting leads to compaction and poor aeration.
Once repotted, trees must be kept frost\--free for at least six weeks while the new root hairs establish.
A ``cold frame'' or an unheated greenhouse is essential during this vulnerable recovery period.

\medskip
\topic{Pruning and the Risk of Bleeding}
\noindent
Mid\--February is your final opportunity for heavy structural pruning on \textbf{Maples (Acer)} and \textbf{Birch (Betula)}.
Once the sap is in full flow, these species will ``bleed'' profusely from large wounds, which can weaken the tree and invite fungal infection.
If you must make a late cut, use \textbf{top\--quality wound paste} immediately.
For other deciduous trees, focus on thinning out the ``crow's feet'' (multiple shoots emerging from one point) to maintain delicate ramification.
Check any wire applied last autumn; as branches thicken with the spring sap rise, wire bite can occur within a matter of days.

\medskip
\topic{Watering and Early Feeding}
\noindent
Watering requirements are increasing.
While we remain in a damp cycle, a sunny afternoon combined with the drying ``March winds'' can desiccate a shallow pot surprisingly quickly.
Check the soil surface every morning.
Regarding nutrition: \textbf{do not feed repotted trees} for at least a month.
For trees not being repotted this year, you may begin a very low\--nitrogen ``awakening'' feed once you see green tips.
This encourages early root activity without forcing lanky, long\--internode growth that ruins the silhouette of a refined specimen.

\medskip
\topic{Managing Desiccating Winds}
\noindent
Late February and early March in our region are often characterised by harsh, dry easterly winds.
These are frequently more dangerous than the frost itself.
A tree with swelling buds has no leaves to lose moisture, but the wind can pull water directly through the bark and bud scales, leading to ``winter kill'' on the windward side of the tree.
If you cannot move your benches, use \textbf{windbreak netting} or move your most prized specimens to the lee of a wall.
Maintain a ``damp but not sodden'' environment; high ambient humidity is your best defence against desiccation during this breezy transition.

\medskip
\topic{Pest and Disease Vigilance}
\noindent
As the temperature rises above 8°C, overwintering pests begin to stir.
\textbf{Aphids} are often attracted to the soft, succulent growth of newly opening buds.
Check the undersides of emerging leaves on Elms and Maples.
If you haven't applied a \textbf{winter wash} or dormant oil spray, do so now before the leaves are fully out, as the oil can scorch tender new foliage.
Keep an eye out for \textbf{mildew} in stagnant air; if using a greenhouse, ensure you open the vents on sunny days to maintain air circulation.
A clean start now prevents a systemic infestation in May.
\closearticle
\columnbreak

